\chapter*{Conferencia VIII}
\addcontentsline{toc}{chapter}{LECTURA I - El timbre y el carácter del violín}
Señores: Nos acercamos a problemas en el timbre del violín que no solo son difíciles de resolver, sino también de enunciar. Durante muchos años permanecí en silencio en mi banco, pensando, trabajando, pero sin escribir. Durante esos años, solo tuve el privilegio, a intervalos largos y breves, de contar con la compañía de personas interesadas en la filosofía del timbre del violín. Como es bien sabido, conversar sobre cualquier tema facilita enormemente la elección de las palabras precisas para que nuestro significado se exprese con claridad. Al carecer de tal beneficio, y al no obtener prácticamente ningún beneficio directo de los libros de filosofía ni de la literatura general, encuentro que la tarea de la dicción precisa es una tarea difícil. Mi edad y mis achaques también dificultan mi pluma. Por lo tanto, si fallo en la precisión, no me sorprenderá. Me reconforta saber que cada nueva generación se vuelve más brillante intelectualmente y, por ende, más capaz de enunciar con precisión tanto los principios como las soluciones. Dado que los físicos han abandonado el problema de la curvatura de las tapas del violín para asegurar la máxima concentración del movimiento molecular en las salidas como irresoluble, de ninguna manera me limito a decir que este problema nunca se resolverá. Creo firmemente que se encontrará una solución. Negar tal posibilidad es tan absurdo como decir que el violín alcanzó la perfección hace 200 años. Según mi experiencia, hoy en día se fabrican violines con mayor potencia tonal y mayor
130

uniformidad de tono mayor que en cualquier otro período del violín.
historia.
[Galilleo no descubrió América, pero su enunciación de principios hizo posible su descubrimiento. Después de que Colón demostrara que los principios de Galileo eran correctos, se cuenta que un grupo de comensales, algo aturdidos, se levantó para comentar que este descubrimiento era algo fácil que cualquiera podría haber hecho; no tenía nada que ver con navegar hacia el oeste hasta encontrar tierra firme y luego regresar. Colón pidió un huevo y les pidió a todos los presentes que lo pusieran de pie sobre su extremo más pequeño. Después de que todos fracasaran, Colón, golpeando el huevo contra la mesa con la fuerza suficiente para romper parcialmente la cáscara, logró que el huevo se mantuviera firmemente de pie sobre su extremo más pequeño. Esta historia mítica se cuenta aquí con un propósito.
Si logro siquiera descifrar la complejidad de algún problema difícil en el sonido del violín, entonces a quien posteriormente consiga que dicho problema se "mantenga en pie" lo llamaré Colón, si tan solo insinúa que mi nombre es Galileo.
Es bueno sonreír mientras tengamos la oportunidad, porque pronto entraremos en un territorio donde la sequedad impide la posibilidad de sonreír. Después de cruzar este "distrito seco" intentaremos encontrar humedad para nuestros labios resecos y celebrar debidamente nuestra "pascua". La uniformidad en los valores tonales del violín se considera uno de los problemas irresolubles. Podemos decir con verdad que este problema ha desafiado con éxito la solución durante 400 años. Que este problema se resuelva algún día, solo podemos esperar. Al intentar una solución...
131

En cuanto a la solución de este problema, logré algo, pero no lo suficiente. Solo en un número limitado de violines he conseguido una uniformidad de timbres que me impide distinguir el tono de cada uno al tocarlo sin que lo vea. En este sentido, el timbre del violín es como la voz de los cantantes. Solo en raras ocasiones oímos a dos cantantes con cualidades vocales idénticas.
¡Qué valor incalculable tiene el violín en la solución de este problema!
Es un valor que desafía las leyes de la aritmética.
Convencido de que este problema interesa a los estudiantes de violín, describo en detalle mi método para lograr la uniformidad en los tonos. En primer lugar, selecciono violines con dimensiones lo más uniformes posible. Es evidente que una gran falta de uniformidad en la capacidad cúbica dificulta la uniformidad de las cualidades tonales. Por lo tanto, cuando dos violines varían mucho en la longitud de las columnas de aire perpendiculares dentro de la caja de resonancia, resulta difícil lograr una afinación uniforme.
[No me refiero a la dificultad de afinar las cuerdas al "unísono", sino a la dificultad de evitar que el oído distinga la presencia de dos sonidos con diferente tono.]
En segundo lugar, selecciono violines cuya madera tenga una madurez, densidad y anchura de veta lo más uniformes posible. Estos últimos puntos influyen mucho en la calidad del sonido del violín. Después de seleccionar así media docena de violines, les quito el barniz. Esta eliminación es necesaria porque tanto la calidad como la cantidad de barniz influyen en la calidad del sonido.
132

para elevar el tono y disminuir el volumen del tono. A continuación, se abren estos violines y se aplica el trabajo necesario a sus superficies interiores. Este trabajo puede variar en cada violín. Algunos pueden necesitar un nuevo enbarrado; otros, una nueva gradación; otros, tanto un nuevo enbarrado como una nueva gradación; posiblemente algunos necesiten forros y bloques de esquina más gruesos; y, en general, la superficie interior se alisa de forma perfecta y permanente mediante un barniz elástico aplicado con una atenuación tal que solo es posible para el
'
El proceso de frotamiento se describió anteriormente. Por supuesto, en el reensamblaje y la regradación, se deben considerar los diferentes grados de densidad de la fibra de la tapa armónica. Después del ensamblaje, se ajusta el diapasón. Su superficie inferior hueca, que comienza en la base del mástil, se deja tan recta que la regla la toque en todos los puntos a lo largo de la cavidad.
[La filosofía que subyace a esta forma de dar forma a la superficie inferior del diapasón aparecerá más adelante en la discusión sobre los modificadores de tono.]
La superficie inferior del diapasón se coloca a una altura uniforme en los seis instrumentos. El peso y la densidad de los diapasones son lo más uniformes posible. Mientras estos violines aún están en bruto, se ajustan puentes de igual madurez, densidad y ancho de veta, y cuerdas de calibre 2, de seis hebras, de torsión dura, cada una de diámetro uniforme, y se realiza la prueba preliminar de uniformidad en los valores tonales. Esta prueba se realiza antes de barnizar la superficie exterior con el fin de conservar la oportunidad de trabajar.
133

sobre dicha superficie para bajar el tono donde sea necesario. Así, después de aplicar el arco sobre cada violín, selecciono uno que tenga un tono promedio, ni el más alto ni el más bajo del lote, y luego subo el tono de los que están por debajo y bajo el de los que están por encima de esta selección. En la tarea de subir el tono, ahora tengo tres, o, contando el barniz, cuatro medios de ayuda, de la siguiente manera:
(1) Ampliar las salidas.
(2) Posición del puente.
(3) Posición del puesto.
(4) Cantidad de barniz.
Para el trabajo de bajar el tono, tengo tres medios, por lo tanto:
(1) Disminución del grosor de la caja de resonancia.
(2) Posición del puente.
(3) Posición del puesto.
[De hecho, un quinto método para elevar el tono puede hallarse contando la disminución de la profundidad de las costillas; (Véase la Regla VI) pero, al ser un trabajo innecesario salvo en casos extremos, no me veo obligado a emplearlo aquí.]
La filosofía involucrada en este trabajo de bajar el tono se reserva para una fecha posterior; por lo tanto, en este momento basta con decir que después de que el tono de estos violines se vuelve tan casi idéntico que mi sentido del oído no puede distinguir el tono de uno de otro cuando se tocan ocultos a la vista, entonces se desmontan y se aplica barniz, de igual calidad y cantidad, sobre la superficie exterior y mediante el proceso de "frotamiento". Cuando están nuevamente listos para usarse,
134

Se les lleva al campo abierto para la prueba de intensidad descrita anteriormente. Antes de presentar los resultados de la prueba de intensidad a larga distancia (capacidad de proyección), mencionaré que en las cualidades de tono e intensidad se observa la mayor proximidad a la igualdad de valores tonales en estos violines, tal como están preparados. A distancias cercanas, aún se aprecia una diferencia notable entre los armónicos graves y los armónicos consonantes. Por lo tanto, su timbre varía en riqueza. Asimismo, existe una diferencia en la respuesta a la presión del arco, particularmente observable en la mayor o menor nitidez de los tonos tocados rápidamente.
Ahora retomo la prueba de larga distancia al aire libre. Antes del tratamiento aquí descrito, el registro muestra que la intensidad del tono en estos violines se dispersa entre 400 y 800 pies. Después de la preparación, como se describe, el registro muestra un aumento en la intensidad del tono hasta 1100 y 1200 pies, con un promedio de 1175 pies, un aumento del 80 por ciento en la distancia; y un aumento, en la uniformidad del valor del tono sobre esta calidad de tono aislada, del 62 por ciento. No conozco ninguna manera de describir con precisión el
Porcentaje de incremento para cualidades tonales como volumen, uniformidad, riqueza, ausencia de armónicos disonantes, simpatía en concierto, respuesta al arco, agradable sonoridad de los acordes dobles, brillantez y potencia armónica. En general, y en mi opinión, estas últimas cualidades tonales aumentan de forma proporcional a la intensidad, y en la mayoría de los casos, creo que el porcentaje de incremento es mayor que el incremento de la intensidad.
135
\label{efecto_tonal}
Hay otro valor tonal en este lote de
Violines. Es un valor difícil de describir con precisión. Me refiero al efecto que se produce al tocar estos violines al unísono. Creo que mi experiencia al escuchar violines preparados de esta manera no tiene parangón. Esta convicción me lleva a dudar antes de intentar describir dicho efecto. En esta descripción, no encuentro apoyo mediante la comparación. Me veo obligado a buscar palabras que describan una nueva impresión tanto en el sentido del oído como en el sentido de la sensibilidad.
La primera palabra que emerge de la niebla es solidez. Cuando decimos que un objeto es "sólido", nos referimos a que posee peso. Por regla general, la palabra sólido se usa solo para describir objetos visibles. Sin embargo, me inclino a usarla para describir el tono combinado de esta media docena de violines. Es cierto que, mientras los escuchaba, a una distancia de más de 300 metros, sentí como si las ondas sonoras me golpearan. Como recordarán, la prueba de larga distancia se realiza al aire libre, bajo un cielo despejado, sin viento, con una temperatura de entre 21 y 32 grados Celsius, con vapor de agua en o por debajo del punto normal (nunca por encima), entre las 10 de la mañana y las 4 de la tarde (horas en las que las ondas sonoras se propagan con mayor dificultad en una fecha determinada), sobre terreno llano y sin objetos circundantes capaces de reflejar las líneas de movimiento molecular. En tales circunstancias, se observa de inmediato que la posición que ocupan estos violines, durante tales
136

La prueba de tono se realiza en el centro de un círculo cuyo radio es mayor de 1000 pies. También se entiende que, bajo las condiciones anteriores, las ondas sonoras de estos violines alcanzan la periferia de este círculo en todos sus puntos.
Soy plenamente consciente, en este momento, de que mi afirmación sobre la «solidez» del sonido de estos violines, preparados según lo aquí expuesto, carece de corroboración. Por lo tanto, les pido que realicen ustedes mismos dicha prueba. En lo que respecta a la aplicación práctica del sonido del violín, mi trabajo ha concluido. Ustedes, conscientes de que en ciertas condiciones las ondas sonoras se propagan con dificultad, y conociendo la queja de no escuchar el sonido del primer violín, retomarán este trabajo donde lo dejo y demostrarán su valor tanto en la orquesta pequeña como en la grande.
En este trabajo, no puedo asegurar que alguien pueda tener éxito sin experiencia. Al reflexionar sobre los años que he dedicado a las particularidades del sonido del violín y observar las dificultades que se presentan en el camino del constructor de violines, me veo obligado a afirmar que el éxito depende en gran medida de la experiencia. Un caballo brioso puede ser espoleado a una acción desbocada por un arnés que le irrita y le queda mal. La caja de resonancia del violín también puede verse afectada por un alma mal ajustada, un diapasón mal ajustado, un puente mal ajustado, un juego de cuerdas inadecuado y un arco de mala calidad.
Cuando todos estos males se combinan, la imaginación falla.
Siempre que una centésima de pulgada de variación en el grosor de la caja de resonancia, en el diámetro de la cuerda, en 137

El grosor y la posición del puente, el grosor y la posición del alma, el área y la posición de las salidas, la profundidad de la costilla y el diámetro de las cerdas del arco influyen en el timbre del violín, mientras el éxito siga siendo un desafío para el maestro artesano. Un carpintero no puede tener éxito en la fabricación de violines, del mismo modo que un herrero no puede tener éxito en la relojería. No quiero decir que un carpintero no pueda convertirse en fabricante de violines, ni que un herrero no pueda convertirse en relojero. Esa es otra cuestión. Pero afirmo categóricamente que ambos deben pasar por un aprendizaje.
Para ilustrar la veracidad de esta última afirmación, basta con mencionar una dificultad común entre los constructores de violines: la sensibilidad variable de la madera de la caja de resonancia. Estas variaciones son tan amplias que pueden causar no solo una profunda sorpresa, sino también una gran decepción. Conozco una sola caja de resonancia, y solo una, que no se ve afectada por cinco veces la cantidad de barniz necesaria para su mera protección. Asimismo, he trabajado con cajas de resonancia tan sensibles que una sola capa de barniz, aplicada con una capa fina, causaba daños evidentes en el sonido. Ante estos hechos, resulta evidente que solo el criterio que emana de la experiencia puede garantizar los mejores resultados. Mi experiencia me lleva a creer que todo buen violín, construido siguiendo reglas estrictas, es fruto de la casualidad.
Es mi deseo ayudar a todos aquellos que intentan el trabajo de regulación del tono del violín. Para este fin, no puedo pensar en una manera más efectiva que señalar hechos como los que ponen en peligro el éxito.
138

Es valioso saber cómo no hacer cierto tipo de trabajo.
En cuanto a la aplicación de material protector a las superficies interiores de violines nuevos, no tengo experiencia que ofrecer. En mi opinión, esto depende del grado de sequedad y de la finalización del proceso de contracción en cada muestra de madera. Evidentemente, no es aconsejable sellar herméticamente la madera del violín antes de que se haya secado y contraído por completo. En nuestro clima, la presencia de vapor de agua en el aire impide tanto la contracción completa como el secado de la madera. Cuando están rodeados de vapor de agua*, los capilares de la madera del violín tienden a absorber y retener más o menos agua. También es evidente que la presencia de agua en los capilares disminuye la resonancia de cada violín. No cabe otra conclusión. Asimismo, es evidente que el sellado hermético de la madera del violín impide eficazmente la entrada de más agua. Por lo tanto, la protección de la superficie interior del violín, además de su valor como medio reflectante perfecto, también es valiosa para mantener la resonancia de forma permanente. Reitero que, antes del sellado hermético, la madera del violín debe estar completamente seca por medios artificiales. En nuestro clima, no podemos esperar que la madera se seque completamente cuando se deja a la intemperie. La observación de una contracción desigual en las tapas de los violines, después de que salen de las manos del constructor, me lleva a la opinión de que, por una cuestión de certeza en los violines nuevos, las tapas, casi reducidas al grosor correcto, deberían dejarse secar durante uno o más años para completar el proceso.
139

de contracción antes de ser ensamblados.
De ninguna manera pretendo que me entiendan como un consejo para esforzarse ininterrumpidamente en producir un sonido de violín con el máximo volumen e intensidad. La necesidad de violines de "gran sonido" solo existe de forma limitada. El número total de auditorios que generan la necesidad de "gran sonido" es pequeño en comparación. Dado que el objetivo de toda interpretación musical es el placer del oído, se deduce que tanto una calidad de sonido desagradable como un sonido tan débil que resulta inaudible frustran el propósito del esfuerzo musical. Por lo tanto, el error de emplear un violín con "gran sonido" para un estudio, salón o auditorio pequeño solo es comparable al error de emplear un sonido débil para auditorios grandes. En salas pequeñas, el sonido "grande" resulta doloroso; en salas grandes, el sonido débil es decepcionante. Esta proposición tiene dos vertientes: una estética y otra comercial. Es bastante seguro asumir que todos los asistentes a las salas de música esperan disfrutar del sonido del primer violín. También es bastante seguro asumir que tanto un sonido desagradable del primer violín como la incapacidad para oírlo contribuyen a la disminución del público. En este caso, no se puede culpar al mecenas por retirarse. Invierte su dinero esperando disfrutar de la música y solo encuentra decepción.
Es evidente, por ambas partes, que la gran mayoría de los violines se utilizan únicamente en habitaciones relativamente pequeñas; en vista de este hecho, es conveniente construir y regular el tono de la gran mayoría de los violines con un volumen e intensidad mínimos.
140

del tono. Hoy en día, el hábil luthier puede acercarse más que nunca a producir un tono de violín "a pedido".
Presento los siguientes factores como factores que demuestran ser fiables para disminuir tanto el volumen como la intensidad del sonido del violín:
(1) acabado.	Control de la acción de la caja de resonancia por var
(2) ing.	Control de la acción de la caja de resonancia mediante flexión
(3) Control de la acción de la caja de resonancia mediante el grosor de la madera.
(4) Control de la acción de la caja de resonancia mediante el arqueamiento de las placas.
(5) barra.	Control de la acción de la mesa de sonido por el
(6) post.	Control de la acción de la mesa de sonido por el
(7) Control de la acción de la caja de resonancia por el puente.
(8) Control de la acción de la caja de resonancia por el diapasón.
(9) Control de la acción de la caja de resonancia mediante el diámetro de las cuerdas.
(10) Disminución del volumen y la intensidad del tono por el estado de las superficies interiores.
(11) 2 Disminución del volumen y la intensidad del tono por el área y la posición de las salidas.
(12)  Disminución del volumen del tono por la profundidad de las costillas.
Cualquiera de estos factores, por sí solo, es capaz de producir una disminución perceptible tanto del volumen como de
141

intensidad del tono del violín.
( 1 ) El control de la acción de la caja de resonancia mediante barniz, en todos los casos excepto en el aquí mencionado, ha demostrado ser posible y bastante fácil de lograr. La única dificultad que he encontrado al disminuir el volumen y la intensidad del sonido del violín mediante barniz elástico se debió enteramente a las variaciones inherentes en la acción de la madera de la caja de resonancia. Estas variaciones son tan amplias que a menudo impiden la regulación del tono mediante reglas "rígidas e inamovibles". En mi experiencia, el mayor efecto del barniz, en el control de la acción de la caja de resonancia, se ha observado en madera de fibra blanda; y el menor efecto se ha manifestado en madera de fibra densa. Por lo tanto, al emplear una cierta cantidad de barniz en cada violín, su efecto sobre el tono varía según los grados de densidad de la madera de la caja de resonancia.
No he encontrado que ninguna cantidad de barniz, aplicada en la parte posterior, disminuya el tono en absoluto. Por el contrario, he observado un ligero aumento en la potencia del tono con capas gruesas de barniz aplicadas en las tapas traseras cuyo grosor se había reducido hasta el punto de causar debilidad tonal. Sin duda, cualquier cantidad de barniz aumenta la rigidez de las tapas del violín; y, en la caja de resonancia, cualquier cantidad de barniz también disminuye la acción independiente de las fibras contiguas. Por lo tanto, el barniz, en la caja de resonancia, disminuye tanto el volumen como la intensidad del tono. A menudo he observado violines con mayor potencia en tonos simples que en dobles. En todos esos casos, he encontrado un aumento
142

de potencia en el tono de doble cuerda tras la eliminación del barniz de la caja de resonancia. Demostrar que el barniz, en la caja de resonancia del violín, reduce tanto el volumen como la intensidad del sonido es algo fácil de lograr.
En su efecto sobre el timbre del violín, el barniz posee un aspecto estético. Muchos músicos prefieren el timbre natural de la madera sin tratar. Cuando dicho timbre posee una calidad excepcional, nadie puede objetar. La belleza de ese sonido es imposible de mejorar por ningún medio conocido. Pero aquí reside la dificultad. La naturaleza caprichosa de la madera de la caja de resonancia produce timbres que van desde lo más bello hasta lo más áspero. Este aspecto pertenece enteramente al ámbito de la naturaleza. El ser humano solo puede modificarlo. Hay un número limitado de músicos que se complacen en la aspereza del timbre del violín; pero la gran mayoría prefiere que dicha aspereza se modifique con barniz y se conforma con la pérdida de volumen e intensidad a cambio de una calidad de timbre menos agradable.
(2) Doblar la caja de resonancia para colocarla en posición es un factor poderoso para disminuir, no solo el volumen y la intensidad del tono, sino también la duración del mismo. Los violinistas, cuyo fuerte reside en la rapidez de ejecución, a veces se conforman con el tono acortado y sin vida que produce la caja de resonancia doblada. La influencia sobre el volumen, la intensidad y la duración del tono al doblar la caja de resonancia es directamente proporcional al grado de doblado. Para disminuir estas cualidades del tono, mi propio gusto
143

No apruebo la tapa armónica curvada, debido a una cierta calidad de sonido "muerta". Prefiero reducir el área de las salidas, dejando libre la acción de la tapa armónica; así se conserva la viveza del sonido. La diferencia en la calidad del sonido favorece enormemente a este último método.
No he observado que doblar la tapa trasera para colocarla en la posición correcta disminuya en absoluto la potencia del sonido; al contrario, doblar una tapa trasera, que ya es demasiado ligera en madera, aumenta la potencia del sonido. Sin embargo, este aumento no se mantiene indefinidamente. La razón es la siguiente: al abrir el violín con la tapa trasera doblada, y apenas unos años después de que salga de las manos del luthier, encuentro que la curvatura se transfiere a la caja de resonancia.
Este resultado podría esperarse si tan solo lo pensáramos; y los dos factores que causan la pérdida de potencia tonal que acompaña a dicha transferencia de inflexión también podrían hacerse evidentes con la reflexión. El problema es mantener el nivel de pensamiento profundo. Es agotador. Los hechos se descubren más fácilmente tropezando con ellos. Tropezar no cuesta nada más que levantarse y comentar.
(3) El control de la acción de la caja de resonancia mediante el grosor de la madera es un factor de gran importancia. Este factor ha recibido más atención que cualquier otro elemento en la lista de modificadores del tono del violín. Solo el efecto de disminuir el volumen y la intensidad del tono por el grosor de la madera de la caja de resonancia será
recibir atención. [En una ocasión posterior, y en la discusión de
máxima uniformidad de tono, este importante tono-


[Se prestará más atención al modificador.] Es evidente que el grosor del violín
La caja de resonancia debe estar determinada por el diámetro de las cuerdas que se vayan a utilizar. Cuando no se indique lo contrario específicamente, me refiero únicamente al calibre 2.
instrumentos de cuerda.
[Como ya he mencionado, mi experiencia me lleva a creer que la función de la caja de resonancia del violín es originar las ondas sonoras que finalmente llegan al oído como tono musical. Antes del fenómeno del barniz n.° 1, traté la tapa trasera como si también fuera un agente productor de sonido. A este trabajo ahora atribuyo la pérdida de calidad tonal de varios violines. Tras limitar el trabajo de regulación tonal exclusivamente a la caja de resonancia, no encontré ningún fallo atribuible a mi trabajo. Al colocar tapas traseras nuevas y resistentes en violines que habían sufrido graves daños en el sonido por intentos de regulación tonal en dicha tapa, logré, en un grado satisfactorio, restaurar la calidad tonal perdida. No quiero que se entienda que digo que la pérdida de calidad tonal no pueda ser consecuencia de un trabajo erróneo en la caja de resonancia. Al contrario, afirmo que la calidad tonal puede perderse fácilmente por un trabajo erróneo en la caja de resonancia. Lo que sí puedo afirmar con optimismo es que, tras limitar el trabajo de regulación del tono a la caja de resonancia, logré un éxito mucho mayor al asegurar el valor tonal en cualidades como la intensidad, el brillo y la dulzura. El volumen del sonido, a corta distancia, lo pude obtener con certeza trabajando en las placas posteriores.


solo; pero, tal tono invariablemente se quedaba muy corto en la prueba de larga distancia al aire libre; mientras que, con una placa trasera fuerte e inflexible, podía obtener el mismo volumen de tono y una intensidad y brillantez de tono mucho mayores mediante el trabajo dirigido a la caja de resonancia. Este hecho lo he demostrado mediante muchas repeticiones de la prueba de larga distancia al aire libre. Como se indicó anteriormente, esta prueba proporciona evidencia basada en la medición real, en pies lineales, de la intensidad (poder de transmisión) del tono del violín. Es una prueba que resuelve esta cuestión sin lugar a dudas. ii ,
Al disminuir el volumen y la intensidad del tono mediante el grosor de la caja de resonancia, no he encontrado la misma dificultad que al intentar obtener el máximo de esas cualidades tonales. Para lograr el máximo volumen e intensidad del tono, a partir de cierto grosor, se debe proceder con precaución en cuanto a las mediciones y la eliminación de madera; esto se debe a que en cada muestra de madera para caja de resonancia existe un grosor óptimo que produce dicho máximo, y por debajo de este grosor se produce una debilidad del tono. En cambio, para obtener un volumen e intensidad de tono reducidos, el grosor puede variar de 16,64 a 10,64, o incluso menos con madera de fibra densa. Estos grandes grosores invariablemente disminuyen tanto el volumen como la intensidad del tono, pero no en igual medida; la mayor disminución se observa en el volumen. Estos grandes grosores también tienden a elevar el tono.
146

En mi experiencia, he observado que cualquier defecto inherente en la calidad del sonido de la madera de la caja de resonancia se desarrolla y se hace evidente a medida que el grosor de la madera disminuye hasta alcanzar el máximo volumen de sonido. Considero que conocer y tener en cuenta este hecho es valioso tanto para el constructor de violines como para los violinistas que priorizan la calidad del sonido sobre la cantidad. He observado que muchos violinistas comparten esta preferencia. Asimismo, he constatado que, para cualquier ocasión y circunstancia, un violinista solo puede estar bien equipado si dispone de varios violines con distintos grados de volumen e intensidad de sonido.
(4). El arqueamiento dado a las placas del violín desafía al lápiz de ese físico que escribiría su solución. Sin embargo, es fácil demostrar que dicho arqueamiento puede aumentar o disminuir tanto el volumen como la intensidad del tono. No conozco ningún otro elemento en la lista de modificadores del tono del violín que ofrezca mayor interés al estudiante de violín que el arqueamiento de la placa. En la aplicación práctica, no hay disminución del interés desde la placa perfectamente plana hasta cada grado hasta una altura de arco igual a 1 pulgada. Manteniendo otras dimensiones iguales, la potencia del tono aumenta, desde ninguna curvatura, hasta cierta altura de arco; y desde ese punto hasta esa altura igual a 1 pulgada, la potencia del tono disminuye constantemente. Otras dimensiones no permanecen iguales, como aumentar el área de salidas, aumentar el sonido
147

Un mayor grosor de la tabla, junto con un aumento del diámetro de las cuerdas, puede dar lugar a una mayor potencia tonal que en la primera propuesta.
Es evidente que el diámetro de las cuerdas puede controlar de forma segura la altura del arco. En mi observación, utilizando la cuerda de calibre 2 y distribuyendo el arco de manera uniforme desde ambos extremos de las placas hasta la posición del puente, el mayor volumen de tono se alcanza a una altura de arco de f pulgadas, mientras que la mayor intensidad de tono se alcanza a una altura de arco de i pulgadas. Deseo que la afirmación anterior se entienda con el entendimiento de que el área de salidas, en cualquiera de las alturas de arco mencionadas, permanece exactamente igual, y que dicha área no es ni grande ni pequeña, y que el ancho y el largo del cuerpo y la profundidad de las costillas no son mayores que los de "tamaño completo"; además, que la distancia entre las salidas, en sus extremos superiores, es exactamente de 1 y 4 pulgadas. Estas condiciones son absolutamente necesarias.
para determinar la influencia del arqueamiento sobre el volumen y la intensidad del tono.
- Presento este violín como ejemplo de la influencia ejercida
sobre la potencia tonal mediante el arco de la placa.
Observas que las salidas son de área media; y
que su posición no es ni "alta" ni "baja".
Como
Si miras este violín de perfil, observas
que su cintura es de concejal; que
su altura de arco, en la posición de la
puente es igual a una pulgada; que la distribución de la
El arco es igual desde los extremos de las placas hasta el
puente; que la calidad de la madera, el barniz y el trabajo-
148

La técnica es buena. Al aplicar el arco, observas que tiene una potencia tonal igual a la de un pewee; e incluso ese tono es "todo interno".
Al intentar explicar esta gran pérdida de potencia tonal, debido a una altura de arco tan formidable, no nos queda más que una suposición. Conocemos el hecho, pero no podemos afirmarlo con certeza.
Además, aunque sé que una mera conjetura no tiene mucho interés, me arriesgaré a presentar mi hipótesis para explicar este problema tan complejo, siempre y cuando no me exijan una afirmación categórica al respecto. Entiendo perfectamente que una afirmación, sin corroboración, no prueba nada. Por lo tanto, comienzo diciendo lo siguiente: creo que la pérdida de potencia tonal en las tapas de violín de arco alto, manteniendo las demás dimensiones iguales, se debe a que las salidas no están alineadas con la trayectoria de la onda sonora.
"Bueno, ahora llega el desafío."
Si, por algún medio posible, pudiéramos situarnos dentro del cuerpo de este violín mientras el arco hace vibrar las cuerdas, y si nuestra visión fuera lo suficientemente rápida como para seguir la trayectoria de cualquier onda sonora que se origina en la superficie interior de la caja de resonancia, entonces sabríamos por qué dicho movimiento continúa viajando indefinidamente "por dentro"; entonces sabríamos por qué la gran mayoría de esos movimientos no logran salir por las aberturas hechas y previstas, especialmente para la salida de las ondas sonoras.
Puedo afirmar con seguridad que esos "si" han agotado hasta el último grano de plumbago en muchos lápices.
